\documentclass[../../../include/open-logic-section]{subfiles}

\begin{document}

\olfileid{sth}{story}{rus}
\olsection{Russell悖论}

在 \olref[sfr][][]{part} 中，我们探讨了一种朴素的集合论。但根据一种\emph{非常}朴素的观点，集合不过是谓词的外延。这种朴素的想法要求我们接受如下原则：

\begin{defish}
  \emph{朴素概括原则。} 对任意公式 $\phi$，$\Setabs{x}{\phi(x)}$ 都存在。
\end{defish}

尽管这条原则很诱人，但可以证明它是不一致的。我们在 \olref[sfr][set][rus]{sec} 中已经见过这一点，但这个结果如此重要，又如此直截了当，值得逐字重复一遍。

\begin{thm}[Russell悖论]
不存在集合 $R = \Setabs{x}{x \notin x}$
\end{thm}

\begin{proof}
如果 $R = \Setabs{x}{x \notin x}$ 存在，那么 $R \in R$ 当且仅当 $R \notin R$，这是一个矛盾。
\end{proof}

Russell在 1901 年 6 月发现了这个结果。（不过，他并没有以我们刚才呈现的那种精确形式提出该悖论，因为他当时考虑的是 Frege在《Grundgesetze》中概述的集合论。我们将在 \olref[blv]{sec} 中回到这一点。）Russell 于 1902 年 6 月 16 日写信给 Frege，解释了 Frege 系统中的不一致性。关于通信及一些背景，参见 \citet[pp.~124--8]{Heijenoort1967}。

值得强调的是，这个两行证明是\emph{纯粹逻辑}的结果。诚然，我们在记号 $\Setabs{x}{x \notin x}$ 中隐含地使用了一条（非逻辑的？）公理，即外延公理；因为如果存在的话，$\Setabs{x}{\phi(x)}$ 被理解为（由外延性保证）满足 $\phi$ 的元素所组成的\emph{唯一}集合。但我们甚至可以避免哪怕一点点对外延性的依赖，只需将结果陈述如下：\emph{不存在一个集合，其元素恰好是所有不以自身为元素的集合}。而这与集合并没有多大关系。正如 Russell 本人所观察到的，完全相同的推理会导致如下结论：\emph{没有人给恰好那些不给自己刮胡子的人刮胡子}。或者：\emph{没有哪只哈巴狗嗅恰好那些不嗅自己的哈巴狗}。诸如此类。从图式上看，该结果的形式仅仅是：
\[
\lnot \exists x \forall z(Rzx \liff \lnot R zz).
\]
而这不过是一阶逻辑的一个（图式）定理。因此，我们不能仅仅通过对集合论修修补补来避免Russell悖论；它甚至在我们\emph{抵达}集合论之前就出现了。如果我们要使用（经典）一阶逻辑，我们就只能\emph{接受}不存在集合 $R = \Setabs{x}{x\notin x}$ 这一事实。

要点在于：如果你既想接受朴素概括原则，又想\emph{避免}不一致性，你就不能只对\emph{集合论}修修补补。相反，你将不得不彻底改革你的\emph{逻辑}。

当然，基于非经典逻辑的集合论已经被提出来了。但它们至少可以说是非标准的。处理Russell悖论的标准方法是将其视为一个直截了当的不存在性证明，然后学着与之共存。这也是我们将采取的做法。

\end{document}